\subsection{The unringed elliptic modular topos}

\begin{result}{E.T7}{Proved}{A direct-product computation retaining all elliptic isotropy}
Put $k=\overline\Q$, use the coordinate $t=j/1728$, and
let $G=\operatorname{Gal}(k(t)^{\mathrm{sep}}/k(t))$, with
the inverse system and inertia classes $[I_x]$ of E.T1.
Define the explicitly specified subgroup
\[
 \begin{split}
 A=\{[\beta]\in\operatorname{Out}_{\mathrm{cont}}(G):\;&
 \beta\text{ permutes }\{[I_x]:x\in k\},\\
 &\beta[I_0]=[I_0],\quad\beta[I_1]=[I_1]\}.
 \end{split}
 \tag{E.T15}
\]
There is an isomorphism
\[
 \Aut\bigl((\mathcal M_{1,1,k})_{\mathrm{et}}\bigr)
                \simeq C_2\times A.                         \tag{E.T16}
\]
Automorphisms of this unringed small etale topos are taken
up to natural isomorphism.  The chosen generic Weierstrass
family below specifies the isomorphism.  The algebraic
group $C_2$ of E.T4 maps isomorphically to its first
factor.  The second factor contains an infinite cyclic
subgroup acting trivially on the completed orbifold
fundamental group of $\mathcal M_{1,1}$.  Thus the local
groups of orders two, four and six, together with all
closed-point specializations, are retained in this
computation.
\end{result}

\begin{proof}
Over $t\ne0,1$ choose the elliptic family
\[
 y^2=x^3-\frac{3t}{t-1}x-\frac{2t}{t-1},
 \qquad \Delta=-\frac{108t^2}{(t-1)^3}.
 \tag{E.T17}
\]
Its $j$-invariant is $1728t$.  It neutralizes the generic
order-two gerbe, so the fundamental group of the generic
topos is $G'=C_2\times G$.  Write the central coordinate
additively.  Choose prolongations and peripheral generators
$c_x$ for every $x\in k$.  The full local decomposition
group is
$D_x=C_2\times I_x$.  The subgroup belonging to the
punctured local uniformizing chart is
\[
 \begin{array}{c|c|c}
 x&K_x\subset D_x&D_x/K_x\rule{0pt}{2.4ex}\\\hline
 x\ne0,1&\overline{\langle(0,c_x)\rangle}&C_2\\
 0&\overline{\langle(1,c_0^3)\rangle}&C_6\\
 1&\overline{\langle(1,c_1^2)\rangle}&C_4.
 \end{array}
 \tag{E.T18}
\]
These maps follow from actual etale uniformizing charts.
At zero set $t=w^3$, quadratically twist (E.T17) by
$w$, and then set $x=w^2X$, $y=w^3Y$.  The result is
\[
 Y^2=X^3-\frac{3w}{w^3-1}X-\frac2{w^3-1}.
\]
It is smooth at zero; its linear coefficient is a
local parameter for the deformation with constant
coefficient nonzero.  Thus $w$ is an etale chart
parameter for the stack.  Its punctured-chart generator
maps to $(1,c_0^3)$: the exponent three comes from
$t=w^3$ and the central coordinate one from the
quadratic twist by $w$.  At one put $t-1=v^2$ and
twist by $v$.  This gives the smooth chart
\[
 Y^2=X^3-3(1+v^2)X-2v(1+v^2).
\]
Its constant coefficient is a local parameter, and its
punctured generator maps to $(1,c_1^2)$.
The two quotient maps of (E.T18) are consequently
$(e,n)\mapsto3e+n$ modulo six and
$(e,n)\mapsto2e+n$ modulo four, after the peripheral
orientations are chosen.  At every other finite point
(E.T17) is a smooth elliptic model, giving the first
row.  These computations identify the actual chart
subgroups and both central twists.

We need the normalizers of these groups.  For a primitive
generator $c$ of a finite-rank free profinite group $F$
and a positive integer $r$, the subgroup
$N=\ker(F\to\Z/r)$ sending $c$ to one and the other
basis elements to zero has a free profinite basis
containing $c^r$.  This follows by completing the usual
Schreier basis; the subgroup has finite index, so its
full profinite topology is induced from the free group.
Corollary~2.3 of \cite{BuxErshovRapinchuk2011} gives
$C_N(c^r)=\overline{\langle c^r\rangle}$.
Multiplying an element of $C_F(c^r)$ by a suitable
power of $c$ puts it in $N$, so
$C_F(c^r)=\overline{\langle c\rangle}$.
A normalizer of $\overline{\langle c^r\rangle}$ acts
by a unit power.  Abelianization, and the injectivity
of multiplication by $r$ on $\widehat\Z$, force that
unit to be one.  Thus its normalizer is the same group.
Using the cofinal finite-puncture quotients as in E.T1
proves $N_G(I_x^r)=I_x$.  It follows from (E.T18) that
\[
                 N_{G'}(K_x)=D_x.                          \tag{E.T19}
\]

We give the reconstruction needed to read the topos group
from these data.  A finite constructible sheaf has a
finite generic $G'$-set $E$, a finite $H_x=D_x/K_x$-set
$B_x$ at each closed point, and $H_x$-equivariant maps
\[
                       B_x\longrightarrow E^{K_x}.
 \tag{E.T20}
\]
Outside a finite set the corresponding inertia is trivial
on $E$ and the map is an isomorphism.  On a local
uniformizing disk the punctured-chart fundamental group
is $K_x$, so the closed stalk of the open direct image
is exactly $E^{K_x}$ with its residual $H_x$ action.
The same open/closed fibre-product gluing as in E.T1,
now equivariant under the finite chart group, proves
this description on objects and morphisms.  Finite
constructible sheaves are the coherent objects, so
these gluing data reconstruct the whole topos.

An equivalence preserves the generic localization and
permutes the closed points.  Their closed topoi have
groups $C_2,C_6,C_4$, so it fixes $0$ and $1$ separately.
The fixed-point functor in (E.T20) is represented by
$G'/K_x$, with its right $D_x/K_x$ action.  Pro-Yoneda
therefore identifies an equivalence with an outer
automorphism of $G'$ permuting the conjugacy classes
$[K_x]$ and their indicated quotient actions.
Conversely such an automorphism transports every datum
(E.T20) and gives an inverse equivalence.
There is no extra kernel: the possible change of a
gluing isomorphism is the normalizer (E.T19), and its
quotient $D_x/K_x=H_x$ is exactly the natural
isomorphism ambiguity on the closed point topos.
Thus the preceding correspondence is an isomorphism
of groups of equivalence classes.

The center of $G$ is trivial, since it maps to the
trivial center of each cofinal free profinite quotient.
Consequently $Z(G')=C_2$ and every automorphism has the
form
\[
 (e,g)\longmapsto(e+\chi(g),\beta(g)),
 \qquad \chi\in\Hom_{\mathrm{cont}}(G,C_2),
 \quad \beta\in\Aut_{\mathrm{cont}}(G).
 \tag{E.T21}
\]
Projection of (E.T18) shows that $[\beta]$ lies in $A$.
Write its action on each inertia generator, after
conjugation, as $c_x\mapsto c_{s(x)}^{\lambda_x}$,
where $\lambda_x\in\widehat\Z^*$ is odd.  For an
ordinary point, preservation of $K_x$ requires
$\chi(c_x)=0$.  At zero it requires
$1+3\chi(c_0)=1$ in $C_2$, hence $\chi(c_0)=0$.
At one, the factor two in $K_1$ imposes no additional
condition.  The Kummer description
\[
 \Hom_{\mathrm{cont}}(G,C_2)=k(t)^*/k(t)^{*2}
\]
shows that exactly two characters survive: every rational
function factors into linear factors over $k$, and a
character unramified at every finite point except one
is represented by $1$ or $t-1$.  Let $\chi_1$ be the
second.  Its values on all finite inertia groups determine
it uniquely.  Since $\beta$ fixes $[I_1]$ and permutes
the other finite inertia classes, $\chi_1\beta=\chi_1$.
Every $\beta\in A$ with $\chi=0$ preserves (E.T18),
and inner changes in $\beta$ lift to inner changes in
$G'$.  Formula (E.T21) therefore gives the direct
product (E.T16), with both inverse maps proved.

The algebraic involution (E.T8), applied to (E.T17),
is its quadratic twist by $\Delta$.
The square class of $\Delta$ is $t-1$, since its
constant is a square in $k$ and $t^2$ is a square.
Thus its image is precisely $(\chi_1,1)$ in (E.T16).
Finally use the rational piecewise-affine homeomorphism
$h$ of E.T2 in the $t$-plane.  It fixes neighborhoods
of $0,1,\infty$, permutes all algebraic points, and
is isotopic to the identity relative to those
neighborhoods.  The same all-puncture algebraization
puts its action in $A$.  Its distinct powers move
subterminal opens, whereas its relative isotopy acts
trivially on the orbifold with its two special points
and generic $C_2$.  Hence it gives the asserted
infinite cyclic subgroup in the kernel of restriction
to finite locally constant sheaves.
\end{proof}
